% Generated draft: requires mathematical and source review.
\documentclass[11pt]{article}
\usepackage[margin=25mm]{geometry}
\usepackage{fontspec}
\setmainfont{Latin Modern Roman}
\usepackage{amsmath,amssymb,mathtools,unicode-math}
\setmathfont{Latin Modern Math}
\usepackage{enumitem,needspace,xurl,hyperref}
\setlength{\parindent}{0pt}
\setlength{\parskip}{5pt}
\setlength{\emergencystretch}{3em}
\newcommand{\N}{\mathbb{N}}
\newcommand{\Z}{\mathbb{Z}}
\newcommand{\Q}{\mathbb{Q}}
\newcommand{\R}{\mathbb{R}}
\newcommand{\C}{\mathbb{C}}
\newcommand{\F}{\mathbb{F}}
\newcommand{\A}{\mathbb{A}}
\newcommand{\PP}{\mathbb{P}}
\newcommand{\E}{\mathbb{E}}
\newcommand{\eps}{\varepsilon}
\newcommand{\dd}{\,\mathrm{d}}
\begin{document}
\section*{UnsolvedMath ID: 10800011}
\textbf{Source problem code:} catalogue entry
\subsection{Definitions and mathematical statement}
\textbf{Problem (Victor Vassiliev, Problem 5A).} Is every hypersurface in $P(d;N)$ connected to a trivial hypersurface by a generic path in $P(d;N)$ along which all changes are Morse surgeries decreasing the source's complexity measure?

Here $P(d;N)$ denotes the hypersurface parameter space used in the article; the supplied excerpt does not define it. Let $C$ denote the source's complexity measure and let $\mathcal T(d;N)$ denote its class of trivial hypersurfaces. These are names for missing source definitions, not newly chosen substitutes. The proposed assertion has the following quantifier structure, for each pair $(d,N)$ in the source's range:
\[
 \forall H\in P(d;N)\quad \exists\gamma:[0,1]\longrightarrow P(d;N),\qquad
 \gamma(0)=H,\quad\gamma(1)\in\mathcal T(d;N),
\]
where $\gamma$ is generic and, at every surgery time $t_i$, the only permitted singular transition is a Morse surgery and
\[
 C\bigl(\gamma(t_i+\eps)\bigr)<C\bigl(\gamma(t_i-\eps)\bigr)
\]
for sufficiently small positive $\eps$ on the two regular sides of that transition. A Morse surgery here has the usual meaning of a topology change across an isolated nondegenerate critical event; the particular allowed indices and local conventions must follow the source's hypersurface setting. The demand is for a wholly decreasing path, rather than merely a path with lower endpoint complexity.

A definitive assertion requires the omitted definitions of $P(d;N)$, $C$, ``trivial,'' and generic paths, including the treatment of singular initial hypersurfaces if they are part of $P(d;N)$. The summary places the question in real hypersurface theory and discusses $N=2$, $d\le4$, and partly $d=6$, but does not decide the ambient affine/projective convention or give the complexity formula. No Betti-number formula or other replacement for $C$ is imposed here.
\subsection{Short English statement}
Can every hypersurface in the specified degree and dimension be simplified to a trivial one using only generic Morse surgeries that strictly lower the prescribed complexity at every step? The parameter space, complexity, and trivial targets are not defined in the extracted excerpt.
\subsection{Sources}
\begin{itemize}
\item Local UnsolvedMath registry, permanent ID 10800011, source code catalogue entry; proposer Victor Vassiliev, no proposal year supplied.
\item Victor Vassiliev, \emph{A Few Problems on Monodromy and Discriminants}, Problem 5A; \url{https://amj.math.stonybrook.edu/html-articles/Files-2015-2024/15-11/}. The registry records HTML-structure extraction and access on 2026-07-29; the article has eighteen labeled theorem-style problems.
\item The supplied research summary attributes relevant low-degree component classifications to Vassiliev's degree-four Morse-polynomial isotopy classification and ``the 64 sextics,'' without furnishing bibliographic references in this record.
\end{itemize}
\end{document}
